\begin{table}[h!]
    \centering
    \begin{tabularx}{\textwidth}{XXXX}
    \hline
        \textbf{English}     &
        \textbf{Sranantongo} &
        \textbf{Ndyuka}      &
        \textbf{Saramaccan}  \\
    \hline
        Now think about a woman who has ten valuable coins. She may lose one of them. What does she do then? She lights a lamp and she sweeps inside her house. She looks carefully until she finds the coin. &
        Efu wan uma abi tin solfrumoni, dan a lasi wan fu den, san a o du? A o leti wan lampu, sibi a oso, èn a o suku bun te leki a feni en. &
        We soseefi, efu wan uman abi tin solufu moni, ne wan kai lasi a ini en osu. Na leti a o leti faya e siibi a hii osu, e suku fu te a fende en baka, no? &
        Nöö söseei tu, ee wan mujëë abu teni kpëngëlë möni nöö hën wan u de lasi, nöö an o sëndë faja nöö a ta bai di wosu köniköni ta suku ën teefa a feni ën baka nöö? \\
        Then she speaks to all her friends and to those that live near to her. She says to them, `I have found the coin that I lost. So come to my house and we can all be happy together.' &
        Te a feni en, a o kari den kompe nanga den birtisma fu den kon, dan a o taigi den taki: `Un prisiri nanga mi, bika mi feni a solfrumoni fu mi di ben lasi.' &
        Da, te a fenden en, da a o kai den mati fi en, da a taigi taki: `U meke piisii anga mi. Bika wan solfu moni di ba lasi de, ma mi fende en baka.' &
        Nöö te a feni ën kaa, nöö a o kai dee maiti fëën ku dee sëmbë dë nëën bandja tuu ko. Nöö a o taki da de taa: `Dee sëmbë o, mi bi abi teni kpëngëlë möni nöö wan u de bi lasi, hën mi ta suku ën tefa u dë aki hën mi feni ën. Nöö un ko wai mbei piizii ku mi makandi.'\\
    \hline
    \end{tabularx}\caption{Luke 15:8-9 in English, Sranantongo, Ndyuka, and Saramaccan.}
\end{table}
